\subsection{Generic two-slice stability near inhomogeneous Cauchy data}
\label{sec:generic-two}

\begin{theorem}
\label{th:generic-two}
Fix any smooth periodic vacuum Gowdy Cauchy background
at a regular time $a>0$ with $Q(a,\theta)$ nonconstant.
There is a common short rational delay $d=\pi/L$ and
a smooth analytic path of admissible Cauchy data starting
at this background, keeping $Q(a)$ fixed, such that
the full orbit snapshot at $a+d$ has a locally
bi-Lipschitz missing-$P_t$ inverse at every point of
the path outside a finite exceptional set.
The velocity and snapshot norms are respectively
$H^s$ and $H^{s+1}$ for integer $s\ge1$.
In particular the given background can be approximated
by admissible backgrounds with stable two-slice
recovery, using the same delay. Stability persists
in an open neighborhood of each nonexceptional point.
\end{theorem}

\begin{proof}
Choose a nonvertical target geodesic whose $Q$ range
contains the range of the given $Q(a)$ on a compact
subarc away from its turning point. This is possible by
taking a sufficiently large upper-half-plane semicircle
with its center to one side of the prescribed range.
Set the comparison initial $P$ to the value on that
arc and take both comparison velocities zero.
These data have a geodesic-image evolution.
They meet the hypotheses of Proposition~\ref{prop:two-slice-open}.

Interpolate the initial $P,P_t,Q_t$ linearly in a
parameter $\zeta\in[0,1]$, from the given data to this
comparison datum, keeping $Q(a)$ fixed.
To impose the vacuum periodic momentum constraint,
replace the trial $Q_t$ by
\[
 Q_t-\frac{\int_{\mathbb T}
        (P_\theta P_t+e^{2P}Q_\theta Q_t)\,d\theta}
       {\int_{\mathbb T}e^{2P}Q_\theta^2\,d\theta}\,Q_\theta .
\]
The denominator is strictly positive on the whole
path. The correction vanishes at the endpoints.
All Cauchy data vary real analytically in $\zeta$ in
all fixed smooth Sobolev norms.

There is a common short time interval of forward
existence, including a small real extension of the
parameter path beyond both endpoints. On a fixed
arc $|\omega_\theta(a)|=e^{P(a)}|Q_\theta(a)|$
has a uniform positive lower bound, since $Q(a)$
is fixed. The expansion \eqref{eq:relative-expansion}
therefore gives \eqref{eq:rotation-cover} uniformly
for a sufficiently short $d>0$ on the path.
Choose $d=\pi/L$ small enough for the same conclusion
and the scalar tangent estimate at the comparison datum.

Let $\mathcal A_\zeta:H^s\to (H^{s+1})^2$
be the derivative of the full terminal orbit snapshot
with respect to the missing initial $P_t$.
The estimate \eqref{eq:rank-two-compact} holds for
all real $\zeta$ in a neighborhood of the path.
Thus $\mathcal A_\zeta$ has closed range and finite
kernel, and its Hilbert normal operator
$B_\zeta=\mathcal A_\zeta^\ast\mathcal A_\zeta$
is Fredholm of index zero.
The first-order energy formulation is real analytic
in Cauchy data on a common small time interval.
For instance, one may complexify the energy spaces and
use the analytic implicit function theorem on the
Duhamel equation on short subintervals, since the
coordinate nonlinearity is analytic on the Sobolev
algebra. This also gives analyticity of the observation
derivative in bounded-operator norm.
The normal operators consequently extend to a holomorphic
Fredholm family on an open connected complex
neighborhood of the real parameter path, using
$B_z=\mathcal A_{\bar z}^\ast\mathcal A_z$.
Fredholmness on a possibly thinner neighborhood follows
from its openness in operator norm.

At the comparison endpoint, the tangent argument in
Proposition~\ref{prop:two-slice-open} gives a
bounded-below $\mathcal A_1$; hence $B_1$ is invertible.
The analytic Fredholm theorem makes the real
noninvertibility set discrete in a neighborhood of
$[0,1]$, and therefore finite on that compact interval.
At any remaining point the derivative is bounded below.
Continuity of the derivative and the same pairwise
integration argument as in
Theorem~\ref{th:fixed-rank-two} give the nonlinear
inverse on its Cauchy neighborhood.
\end{proof}

There is no contradiction between density of stable
nearby backgrounds at a short chosen delay and the exact
different histories at the specially resonant delays.
The theorem varies a Cauchy background at a fixed short
delay; it makes no claim about a specified background
and a specified arbitrary observation gap.
